\documentclass[../main.tex]{subfiles}

\begin{document}
    \begin{center}
        \Large{\textbf{ABSTRACT}}
    \end{center}
    
    \vspace{1cm}
    
    \noindent Threshold batch identity-based encryption provides a natural way to distribute decryption authority among multiple parties while preserving the usability of identity-based public keys as well the efficency . This primitive is especially relevant to blockchain encrypted mempools, where transaction contents should remain hidden until their ordering has already been fixed, thereby reducing the risk of front-running, sandwich attacks, and other forms of Maximal Extractable Value. However, many existing batch threshold encryption schemes rely on pairing-based assumptions, which are vulnerable to quantum attacks, while thresholdizing lattice-based trapdoor primitives remains a challenging problem.

This thesis studies a post-quantum approach to threshold identity-based encryption by applying the thresholdizing technique from \cite{AC:ALLW25} on schemce from \cite{EPRINT:BonLauTas25}.Particularly, we build on recent lattice-based batched decryption techniques and combine them with partial lattice trapdoors, which allow the preimage-sampling capability of a lattice trapdoor to be distributed among several parties. Each authority only holds a partial trapdoor and can locally generate a partial decryption share; only a qualified threshold set can combine these shares to recover the final batch decryption key. This gives a direct and algebraic method for constructing a threshold version of lattice-based batched identity-based encryption.

\end{document}